\documentclass[11pt,letterpaper]{article}
\usepackage[margin=1in]{geometry}
\usepackage{amsmath,amssymb,amsthm}
\usepackage{hyperref}
\usepackage{cleveref}
\usepackage{booktabs}

\numberwithin{equation}{section}
\renewcommand{\theequation}{S15.\arabic{equation}}

\title{Supplement S15: The Higgs Mechanism as Angular Capacity \\
\large Supporting Manuscript \S18}
\author{UDT Collaboration}
\date{}

\begin{document}
\maketitle

\begin{abstract}
We derive the Higgs vacuum expectation value
$v = 504\pi^6 m_e = 247.6$~GeV ($+0.6\%$), the self-coupling
$\lambda = \pi/24$ ($+0.7\%$), and the Higgs boson mass
$m_H = 126.7$~GeV ($+1.1\%$) from the angular quantum numbers.  The
VEV is the total angular capacity of the metric's eigenvalue structure.
The Yukawa couplings are capacity fractions.  Spontaneous symmetry
breaking is the Diophantine selection itself.  Verified by
\texttt{analysis/21\_higgs\_sector.py}.
\end{abstract}

\tableofcontents

%----------------------------------------------------------------------
\section{The Standard Model Higgs}
%----------------------------------------------------------------------

In the Standard Model, the Higgs mechanism introduces:
\begin{itemize}
  \item A vacuum expectation value $v = 246.22$~GeV (measured, not derived).
  \item A quartic self-coupling $\lambda \approx 0.13$ (measured).
  \item Twelve Yukawa couplings $y_f$ (measured, not derived).
  \item A Mexican hat potential $V(\Phi) = -\mu^2|\Phi|^2 + \lambda|\Phi|^4$.
\end{itemize}
The entire electroweak symmetry breaking mechanism --- how the $W$ and
$Z$ bosons acquire mass, how fermions acquire mass --- rests on these
$14+$ free parameters.  The Standard Model explains none of them.

%----------------------------------------------------------------------
\section{The VEV as total angular capacity}
%----------------------------------------------------------------------

The Higgs VEV is the maximum mass the metric cavity's angular structure
can support --- the total capacity of the three-index coupling space:
\begin{equation}\label{eq:vev}
  v = (2j+1)^3(2\ell+1)^2(2|\kappa_\mathrm{max}|+1)\,\pi^6\,m_e
  = 504\,\pi^6\,m_e\,.
\end{equation}

The multiplicity $504$ decomposes as:
\begin{equation}
  504 = 8 \times 9 \times 7
  = (2j+1)^3 \times (2\ell+1)^2 \times (2|\kappa_\mathrm{max}|+1)\,.
\end{equation}

\begin{itemize}
  \item $(2j+1)^3 = 8$: spin-$1/2$ at a three-field vertex
    ($2^3$, three Dirac fields meet).
  \item $(2\ell+1)^2 = 9$: orbital exchange squared ($\ell = 1$,
    the same factor as in $1/\alpha$ and $f_\pi$).
  \item $(2|\kappa_\mathrm{max}|+1) = 7$: the closure orbit
    (the same number as CMB peaks).
\end{itemize}

The power $\pi^6$ is the maximum angular power: the full
$S^2$ integration carried to all six angular coupling channels.

\textbf{Result:} $v = 504\pi^6 \times 0.51100~\text{MeV} = 247.6$~GeV,
compared to $246.22$~GeV ($+0.6\%$).

%----------------------------------------------------------------------
\section{Self-coupling and boson mass}
%----------------------------------------------------------------------

The Higgs self-coupling is:
\begin{equation}\label{eq:lambda}
  \lambda = \frac{\pi}{24}\,,
  \qquad 24 = (2j+1)^3(2\ell+1) = 8 \times 3\,.
\end{equation}
Numerically $\lambda = 0.1309$; the LHC measurement gives
$\approx 0.13$ ($+0.7\%$).

The Higgs boson mass follows from the standard relation:
\begin{equation}\label{eq:mH}
  m_H = v\sqrt{2\lambda} = v\sqrt{\frac{\pi}{12}}
  = 126.7~\text{GeV}\,,
\end{equation}
compared to $125.25$~GeV ($+1.1\%$).

\subsection{The top quark relation}

The ratio of the VEV to the top quark mass:
\begin{equation}
  \frac{v}{m_t} = \frac{(2\ell+1)}{(2j+1)} = \frac{3}{2}\,,
\end{equation}
predicting $m_t = v/(3/2) = 164.1$~GeV.  The experimental value
$172.7$~GeV gives $v/m_t = 1.426$, about $5\%$ from $3/2 = 1.5$.
This reflects the approximate nature of the absolute quark mass
formulas (the \emph{ratios} are sub-$2\%$).

%----------------------------------------------------------------------
\section{Yukawa couplings as capacity fractions}
%----------------------------------------------------------------------

In the Standard Model, the Yukawa coupling of each fermion $f$ is
$y_f = \sqrt{2}\,m_f/v$ --- a free parameter for each of the
12~fermion masses.  In UDT, every $m_f$ is derived from
$(j, \ell, |\kappa_\mathrm{max}|)$ and $m_e$, and $v$ is derived from
the same quantum numbers.  The Yukawa couplings are therefore
\emph{ratios of derived quantities}:
\begin{equation}
  y_f = \frac{\sqrt{2}\,m_f}{v}
  = \frac{\sqrt{2}\,(\text{angular/eigenvalue formula for }m_f)}%
         {504\,\pi^6\,m_e}\,.
\end{equation}

Examples:
\begin{align}
  y_\pi &= \frac{\sqrt{2} \times 84\pi}{504\pi^6}
  = \frac{\sqrt{2}}{6\pi^5}\,, \\
  y_p &= \frac{\sqrt{2} \times 6\pi^5}{504\pi^6}
  = \frac{\sqrt{2}}{84\pi}\,.
\end{align}

The Yukawa couplings are not fundamental.  They are geometric ratios,
all already derived.

%----------------------------------------------------------------------
\section{Symmetry breaking as Diophantine selection}
%----------------------------------------------------------------------

In the Standard Model, spontaneous symmetry breaking is a dynamical
process: the Higgs potential has a continuous family of degenerate
minima (the ``Mexican hat''), and the field randomly selects one.
Nobody explains why that specific minimum or that specific potential
shape.

In UDT, the symmetry breaking is \emph{geometric selection}.  The
Diophantine identity
\begin{equation}
  (2j+1)^2(2\ell+1)(2|\kappa_\mathrm{max}|+1)
  = \binom{2\ell + 2|\kappa_\mathrm{max}| + 1}{2\ell + 1}
\end{equation}
has a \emph{unique} solution for spin-$1/2$ fermions on $S^2$:
$(j, \ell, |\kappa_\mathrm{max}|) = (1/2, 1, 3)$.  That uniqueness
IS the breaking of symmetry.  There is only one consistent angular
structure.  The universe did not randomly pick a vacuum --- it is the
only one that exists.

%----------------------------------------------------------------------
\section{Honesty note}
%----------------------------------------------------------------------

The derivations $v = 504\pi^6 m_e$ and $\lambda = \pi/24$ are
candidates: the multiplicities and $\pi$-powers are consistent with
the quantum number structure, but a derivation from the metric action
principle (showing that these specific combinations minimize an
action functional) is open.  The results are reported at the
candidate level.

%----------------------------------------------------------------------
\section{Verification}
%----------------------------------------------------------------------

All results verified by \texttt{analysis/21\_higgs\_sector.py}.
Output: \texttt{data/generated/21\_higgs\_sector.json}.

\end{document}
